\subsection{重数与平坦扩张}

命题 4. — 设 \(\rho: A \to B\) 是诺特局部环之间的局部同态，且 \(N\) 是一个 \(B\)-模，在 \(A\) 上平坦，并且 \(N \otimes_A \kappa_A\) 是有限长度的 \(B\)-模。若 \(M\) 是一个非零有限生成的 \(A\)-模，且 q 是 \(A\) 的一个不同于 \(A\) 的理想，并且 \(M/qM\) 具有有限长度，则 \((M \otimes_A N)/(qB)(M \otimes_A N)\) 是一个非零有限生成的 \(B\)-模，具有有限长度，并且有

\[
e _ {\mathrm{qB}} ^ {\mathrm{B}} (\mathbf {M} \otimes_ {\mathrm{A}} \mathbf {N}) = \operatorname{long} _ {\mathrm{B}} (\mathbf {N} \otimes_ {\mathrm{A}} \kappa_ {\mathrm{A}}). e _ {\mathrm{q}} ^ {\mathrm{A}} (\mathbf {M}).
\]

设 L 是一个有限长度为 r 的 A-模。则 L 有一个长度为 r 的 Jordan–Hölder 序列，其商同构于 \(\kappa_{A}\)；由于 N 在 A 上平坦，B-模 \(L \otimes_{A} N\) 有一个长度为 r 的合成序列，其商同构于 \(N \otimes_{A} \kappa_{A}\)，因而其长度为 \(r \cdot \operatorname{long}_{\mathbf{B}}(\mathbf{N} \otimes_{\mathbf{A}} \kappa_{\mathbf{A}})\)。由于 B-模 \((\mathbf{M} \otimes_{\mathbf{A}} \mathbf{N})/(\mathbf{q}\mathbf{B})^{n}(\mathbf{M} \otimes_{\mathbf{A}} \mathbf{N})\) 同构于 \((\mathbf{M}/\mathbf{q}^{n}\mathbf{M}) \otimes_{\mathbf{A}} \mathbf{N}\)，且这一同构对每个 \(n \in \mathbf{N}\) 都成立，根据重数的定义，命题得证。

推论。— 假设 B 在 A 上平坦，且 \(\rho(\mathfrak{m}_{\mathbf{A}})\) B = \(\mathfrak{m}_{\mathbf{B}}\)。则

\[
e _ {\mathbf {q B}} ^ {\mathbf {B}} (\mathbf {M} \otimes_ {\mathbf {A}} \mathbf {B}) = e _ {\mathbf {q}} ^ {\mathbf {A}} (\mathbf {M}).
\]

这特别适用于 B 是 A 关于一个不同于 A 的理想的完备化 * 或 Hensel 化 *，或 A 的一个膨胀 *（例如 A 的严格 Hensel 化）时。 *

例。— * 设 X 是复代数簇，\(\mathcal{O}_{X,x}\) 是 X 在有理点 \(x\) 处的局部环，\(X^{an}\) 是与 X 关联的解析空间；仍以 \(x\) 表示 \(X^{an}\) 中对应于前述 \(x\) 的点，并令 \(\mathcal{O}_{X^{an},x}\) 为 \(X^{an}\) 在 \(x\) 处的局部环。则 \(e(\mathcal{O}_{X^{an},x}) = e(\mathcal{O}_{X,x})\)。 *

